\section{Exact canonical recipes determine entire histories and finitely many attained maxima}

\begin{definition}[Literal empty or full-maximum recipe]
\label{hir:canonical-word-history-recipe}
\uses{hir:tower-isomorphism}
\lean{Hironaka.HironakaCanonicalWordHistoryUniqueness.CanonicalWord}
\leanok
For an actual marked record $M$, actual history invariant and actual
finite blowup tower $T$, a Boolean word records the literal recipe at
every step, in reverse chronological order. False prescribes the empty
centre; true prescribes the entire actual invariant maximum on the
current actual cosupport. A true maximum may itself be empty. The
identity tower has the empty word. No historical slot is deleted.
\end{definition}

\begin{theorem}[Actual recipe words and weak canonical histories]
\label{hir:canonical-word-history-properties}
\uses{hir:canonical-word-history-recipe}
\lean{Hironaka.HironakaCanonicalWordHistoryUniqueness.CanonicalWord.isCanonical,Hironaka.HironakaCanonicalWordHistoryUniqueness.CanonicalWord.length_eq,Hironaka.HironakaCanonicalWordHistoryUniqueness.exists_canonicalWord}
\leanok
Every actual recipe gives a weak canonical history, its word length
is exactly the actual tower length, and every weak canonical history
has such a word. Empty and full recipes retain all historical positions.
\end{theorem}
\begin{proof}
Induct on the actual tower. At the identity the empty word gives the
identity weak canonical history and both lengths are zero. At a step,
the preceding word is supplied by induction. The literal recipe says
that the actual centre is empty or the entire current maximum, which
is precisely the next weak canonical condition. Both lengths increase
by one. Conversely the actual weak canonical condition supplies its
empty-or-full alternative; prepend the corresponding Boolean to the
preceding word. This includes the case in which both alternatives
describe an empty actual centre.
\end{proof}

\begin{theorem}[The same recipe determines the entire actual history over a biholomorphism]
\label{hir:canonical-word-history-whole-uniqueness}
\uses{hir:canonical-word-history-properties,hir:tower-isomorphism}
\lean{Hironaka.HironakaCanonicalWordHistoryUniqueness.exists_isTowerIso_of_canonicalWord}
\leanok
For an original marked record in the actual invariant domain and a
genuine base biholomorphism, any two actual histories for the record
and its literal pullback that realize the same word have a genuine
whole tower isomorphism over that base biholomorphism. No nonempty-centre
assumption, selected comparison or terminality is supplied.
\end{theorem}
\begin{proof}
Induct on the common word and both actual towers. The empty word forces
both towers to be identities, whose comparison is the given base map.
For the next bit the induction hypothesis supplies a genuine isomorphism
of the entire preceding histories. Its actual whole-record and cosupport
dictionary, together with the true invariant history-isomorphism law,
identifies the two entire current maximum loci by inverse image. For
a false bit the two actual centres are empty, so inverse-image equality
holds as well. The actual universal uniqueness of blowup along these
corresponding centres constructs a genuine biholomorphism of the next
spaces. Append its true projection commutation and centre equality to
the preceding tower isomorphism. This proves the whole-history result
without compressing any empty positions.
\end{proof}

\begin{theorem}[Only finitely many entire attained maxima occur at each fixed historical length]
\label{hir:canonical-word-finite-attained-maxima}
\uses{hir:canonical-word-history-whole-uniqueness}
\lean{Hironaka.HironakaCanonicalWordHistoryUniqueness.finite_attained_maximum_values}
\leanok
For any original marked record in the actual invariant domain and any
fixed natural length $j$, the set of entire maximum values attained by
any actual weak canonical history of length $j$ is finite. This does
not assert that any particular noncompact history has an attained maximum.
\end{theorem}
\begin{proof}
Every such actual history supplies a Boolean word of exactly length
$j$. For a fixed word, two possible histories are genuinely isomorphic
over the identity by the preceding theorem. Their entire cosupports
and invariant values correspond. Hence any attained maximum value in
one equals any attained maximum value in the other: compare the image
of a maximizing point with a maximizing point of the other history.
The finite set of Boolean words of length $j$ therefore has at most
one attained maximum value per word. Choose one realizing word for
each candidate value. The established uniqueness makes this assignment
injective, proving finiteness of the unchanged set of actual attained
maximum values. No finite value-image or global maximum is assumed.
\end{proof}
